% =====================================================================
\chapter[内部基础设施化]{内部基础设施化：内核六职能的条件必然性}\label{ch:infrastructure}
% =====================================================================

本章处理问题 \ref{prob:scale}。核心主张是：任何足够复杂的软件系统都会在自己的问题域
内重新发明一个操作系统。其原因不是工程师在模仿内核，而是一组语义义务在规模上无法
被回避。本章给出该主张的前提、证明与判据，并给出一个可在单一项目内实施的可证伪检验。

\section{现象与术语}\label{sec:phenomenon}

观察到的现象是：当系统规模增长到一定程度，它会开始把原本分散在各处的机制提升为
统一的系统级能力。具体而言：

\begin{itemize}
\item 多个功能都需要分配和释放资源；
\item 多个子系统都需要管理对象生命周期；
\item 多个组件都需要执行异步任务；
\item 多个模块都需要共享状态并协调修改；
\item 多个功能都需要统一的错误处理、日志与诊断机制；
\item 多个执行路径都需要遵循相同的线程安全与并发规则。
\end{itemize}

如果每个模块独立解决这些问题，系统就会出现大量重复实现、不一致的行为与难以控制的
耦合。因此系统开始将这些问题集中处理，形成公共机制、接口、管理器、运行时与执行模型。

\begin{definition}[内部基础设施化]\label{def:infra}
\term{内部基础设施化}指一个系统依次经历下列三个阶段：
\begin{enumerate}[label=(\roman*)]
\item \term{共性识别}：把重复出现的需求识别为同一类问题（问题空间的归并）；
\item \term{机制统一}：把共性提升为公共能力，使不同功能通过公共接口访问共享资源，
      遵循共同的生命周期、错误处理与并发规则（机制的集中化与标准化）；
\item \term{结构统一}：上层功能不再直接解决全部底层问题，而在统一规则之上构建
      （结构上的统一）。
\end{enumerate}
简记为 $\boxed{\text{共性识别} \to \text{机制统一} \to \text{上层组合}}$。
\end{definition}

机制统一\emph{不}意味着所有逻辑都集中在一个中心模块中。成熟的
架构也可以\emph{通过协议、接口与分布式所有权}实现统一。统一的是\emph{规则与契约}，
不一定是\emph{实现与所有权}。这一区分是避免"上帝对象"的关键，第 \ref{ch:design} 章
会把它写成一条设计原则。

\section{参与者-资源结构}\label{sec:participant}

\begin{definition}[参与者-资源结构]\label{def:pr-structure}
一个\term{参与者-资源结构}是一个五元组
\begin{equation}\label{eq:pr}
\Pi = \bigl(E,\; \mathcal{R},\; \to,\; \mathrm{Ready},\; \mathrm{Fault}\bigr),
\end{equation}
其中 $E$ 是并发存活实体的集合，$\mathcal{R} = \{r_1,\dots,r_m\}$ 是有限共享资源集合，
$\to$ 是操作语义的转移关系，$\mathrm{Ready} \subseteq E \times \mathrm{Conf}$ 指出在
给定配置下哪些实体可执行，$\mathrm{Fault}$ 是故障事件集合。
\end{definition}

关键的两个量是 $n = |E|$ 与 $m = |\mathcal{R}|$。第 \ref{sec:integration} 节将证明：
决定是否需要统一基础设施的自变量是 $n \times m$ 的结构复杂度，而不是代码行数。

\begin{definition}[统一基座]\label{def:substrate}
$\Pi$ 的一个\term{统一基座}是一组机制 $\mathcal{U}$，使得每个实体 $e \in E$ 与每个
资源 $r \in \mathcal{R}$ 的交互都通过 $\mathcal{U}$ 声明的接口进行，而不是由 $e$
直接实现。
\end{definition}

\section{集成成本定理与盈亏平衡阈值}\label{sec:integration}

\begin{theorem}[集成成本]\label{th:integration}
设参与者-资源结构 $\Pi$ 有 $n$ 个实体与 $m$ 个共享资源。假设
\begin{enumerate}[label=(\alph*)]
\item 在无基座情形下，每个实体若参与资源 $r$，必须独立实现与 $r$ 的正确交互，
      包括对 $r$ 的不变量的独立处理；
\item 在最大耦合情形下，每个实体参与全部 $m$ 个资源；
\item 有基座情形下，每个实体至多使用 $k$ 个基座接口，且基座自身对 $m$ 个资源各定义
      一次语义；基座的固定成本为 $F$（设计、实现与长期维护）。
\end{enumerate}
则无基座情形的交互对数为 $\Theta(nm)$，有基座情形为 $\mathcal{O}(nk + m)$。记
$a = \sum_{e} |\mathcal{R}(e)| \le nm$ 为实体-资源关联数，$c$ 为单次直接交互的单位
成本，$c_s$ 为单次基座接口的单位成本，$C_{\mathrm{cons}}$ 为维持一个资源的 $n$ 份独立
不变量实现彼此一致的成本，$F$ 为基座的固定成本，则两种安排的成本为
\begin{equation}\label{eq:costs}
C_{\text{直}} \;=\; c\,a + C_{\mathrm{cons}},
\qquad
C_{\text{基}} \;=\; F + c_s\,(n k + m).
\end{equation}
基座是期望意义上更好的安排，当且仅当
$C_{\text{基}} < C_{\text{直}}$，即当且仅当
\begin{equation}\label{eq:threshold}
c\,a + C_{\mathrm{cons}} \;>\; F + c_s\,(n k + m).
\end{equation}
\end{theorem}

\begin{proof}
无基座情形下交互对数为 $\sum_{e} |\mathcal{R}(e)| = a \le nm$，在假设 (b) 下取
等号；每对需要成本 $c$ 才能被正确实现，此外，对每个资源 $r$，$n$ 份独立实现必须对
$r$ 的不变量达成一致，这耗费 $C_{\mathrm{cons}}$，其失败模式数目为 $\Omega(n)$
（任一实现偏离即破坏全局不变量）。有基座情形下 $m$ 个资源的语义各定义一次，每个
实体通过至多 $k$ 个接口与基座交互，固定成本 $F$ 只付一次，即得式 \eqref{eq:costs}。
比较两式即得式 \eqref{eq:threshold}。
\end{proof}

需要对式 \eqref{eq:threshold} 的形式作出两点说明，因为那些显而易见的简化各自以不同
方式出错。

第一，两边必须是成本，而不是计数。本结果早先的表述把对数 $nm$ 直接与
$c(nk+m) + F$ 相比；这在量纲上不一致，因为左边在数交互次数，右边在为它们定价。正确
的比较是在 $C_{\text{直}}$ 与 $C_{\text{基}}$ 之间，如式 \eqref{eq:threshold} 所示。
若取两个单位成本相等，$c = c_s = c$，则不等式化为
\begin{equation}\label{eq:threshold-equal}
a + \frac{C_{\mathrm{cons}}}{c} \;>\; n k + m + \frac{F}{c},
\end{equation}
而在假设 (b) 下，此时 $a = nm$，
\begin{equation}\label{eq:threshold-canonical}
n m \;>\; n k + m + \frac{F}{c} - \frac{C_{\mathrm{cons}}}{c}.
\end{equation}
右端三项分别是基座的按实体接口开销、按资源开销，以及以一次交互为单位计的固定成本；
最后一项是基座所消除的一致性成本。舍去最后一项会使式 \eqref{eq:threshold-canonical}
偏于保守。

第二，$a$ 不必等于 $nm$。假设 (b) 是最坏情形，而不是典型情形。若实体 $e$ 只参与
$|\mathcal{R}(e)|$ 个资源，且 $\mathbb{E}_e[|\mathcal{R}(e)|]$ 不随 $m$ 增长，则 $a$
按线性而非按乘积增长，基座就完全不必收回自身成本。因此，资源使用是弥散还是集中，是
一个关于系统的经验问题，而不是代数推导的后果；附录 \ref{app:proofs} 记录了乘积形式
得以成立所要求的增长条件。

\begin{corollary}[盈亏平衡阈值]\label{cor:threshold}
由于 $k$ 是每个实体的\emph{有界}接口使用量，$\mathcal{O}(nk + m)$ 关于 $n$ 是线性的，
而 $\Theta(nm)$ 关于 $n$ 与 $m$ 都是线性的乘积。因此在假设 (b) 与单位成本相等的
条件下，只要 $n$ 或 $m$ 之一显著增长，式 \eqref{eq:threshold-canonical} 即被满足。
盈亏平衡点位于
\[
c\,\bigl(a - n k - m\bigr) \;\sim\; F,
\]
即当省下的交互工作量按单位交互成本计价后超过基座的固定成本时。在假设 (b) 下且低阶项
被支配时，此式读作 $n m \sim F/c$，故起支配作用的量是\emph{实体数与资源数的乘积}。
\end{corollary}

该推论被有意命名为盈亏平衡，而不是相变。两条成本函数相交，决定的是一种安排开始变得
更便宜的那一点；而\emph{相变}还额外要求系统的组织在该点的某个邻域内发生结构变化。
这种变化是否发生——越过盈亏平衡点的系统是否真的长出一个基座——是一个经验问题，即
假设 $H_3$；这里的代数给出的是使该问题变得尖锐的那个阈值，而不是问题的答案。

推论 \ref{cor:threshold} 精确化了"代码量达到某个数量级"这一直觉，并给出了它的正确
形式：

\begin{principle}[阈值不是代码行数]\label{pr:notloc}
决定是否必须建立统一基座的自变量是\emph{活跃实体数}与\emph{共享资源数}的乘积，
以及\emph{重复问题的出现频率}、\emph{状态与资源之间的交互数量}、\emph{并发执行的
复杂程度}、\emph{模块的演化频率}与\emph{对一致性、性能与可靠性的要求}。代码行数
只是这些量的代理变量。
\end{principle}

这解释了一个本该令人困惑的事实：两个各有一百万行代码的软件，内部架构可能完全不同。
一个可能是大量彼此独立的生成代码、数据文件或线性业务逻辑（$m$ 小，$n$ 小）；另一个
可能有复杂的共享状态、异步任务与大量互相依赖的子系统（$n$、$m$ 都大）。只有后者会
出现对统一运行时、资源管理与协调机制的需求。

核心归并（\ref{sec:participant}）在这里给出第二个解释：当系统中活跃实体的数量从
$10^{2}$ 跃迁到 $10^{6}$ 时，朴素管理方式发生\emph{相变}。
管理 100 个对象时，条件分支与裸分配是足够的；管理 100 万个并发任务、缓存条目或连接
时，实体间的潜在交互是 $O(N^{2})$，复杂度不可行，于是被迫引入池化、索引、分层、异步、
批量与缓存局部性，而这套手段的集合恰恰构成"内核技术"的全部清单。

\begin{example}\label{ex:pg-memctx}
一个关系数据库引擎自建内存上下文（memory context）体系，是因为通用分配器在"每条
查询创建并销毁数千个元组"的负载下会产生碎片且无法追踪泄漏。于是它构造了上下文树的
分配与批量回收机制，这是数据库域内的对象池。一个异步运行时自建工作窃取调度器，
是因为操作系统线程在百万并发面前过重，于是构造了 $M{:}N$ 调度，这是用户态的分时
调度。数量级跃迁强制引入内核技术，这是第一重必然。
\end{example}

\section{内核六职能}\label{sec:six}

\begin{definition}[内核六职能]\label{def:six}
对参与者-资源结构 $\Pi$，定义\term{内核六职能} $\mathcal{K}_6$ 为
\begin{enumerate}[label=(F\arabic*)]
\item \term{身份与生命周期} \textnormal{(identity \& lifecycle)}：实体的身份标识、
      创建、存活与终止关系；
\item \term{分配与回收} \textnormal{(allocation \& reclamation)}：有限资源的分配
      策略与回收策略；
\item \term{调度} \textnormal{(scheduling)}：在就绪实体之间选择下一个执行者的选择
      函数；
\item \term{命名与寻址} \textnormal{(naming \& addressing)}：使实体可被发现、可被
      引用、可被遍历的机制；
\item \term{通信与同步} \textnormal{(communication \& synchronisation)}：实体之间的
      交互原语与并发转移的复合规则；
\item \term{故障隔离与恢复} \textnormal{(isolation \& recovery)}：故障的包含关系与
      状态的可重放性。
\end{enumerate}
\end{definition}

\begin{theorem}[六职能的条件必然性]\label{th:six}
设 $\Pi = (E, \mathcal{R}, \to, \mathrm{Ready}, \mathrm{Fault})$ 为一个参与者-资源
结构，其操作语义良定义（即在每个可达配置下，下一步转移是被确定的，或其所允许的
非确定性已被显式规定）。若 $|E| > 1$ 且 $\mathcal{R}$ 有限且受争用，则 $\mathcal{K}_6$
的六项职能必然已被定义，定义位置或在该系统自身的代码中，或者在其宿主（操作系统、语言
运行时、框架）中。
\end{theorem}

\begin{proof}[结构论证]
逐项说明缺失任一职能将导致语义不完备（存在未定义行为）。

\emph{(F1) 身份与生命周期}：若实体不可区分，则 $\mathrm{Ready}$ 无法被定义为一个
关于实体集合的函数；若实体无创建与终止关系，则 $E$ 随时间变化而未被规定，$\to$ 在
"实体出现/消失"的步骤上无定义。

\emph{(F2) 分配与回收}：$\mathcal{R}$ 有限且受争用，故"同时请求超过可供给量的资源"
是一个可达状态。若缺少分配策略，$\to$ 在该状态下无定义。终止实体的资源若不回收，
则可用资源单调递减，系统在有限步内进入资源耗尽状态；此时 $\to$ 在前瞻意义下无定义。

\emph{(F3) 调度}：若 $|E| > 1$ 且多个实体同时就绪，则 $\to$ 不再是一个函数而是一个
关系。若不规定就绪者之间的选择，则"允许的非确定性"未被界定，语义不完备。

\emph{(F4) 命名与寻址}：交互原语（见 F5）必须指定交互的对象。若无引用机制，则交互
的宾语不可表达。

\emph{(F5) 通信与同步}：若实体之间存在交互，则全局转移不是局部转移的简单积；缺少
复合规则时，并发转移的语义未被定义。

\emph{(F6) 故障隔离与恢复}：$\mathrm{Fault} \ne \varnothing$（故障是物理与工程的
既成事实）。若缺少包含关系，则一次局部偏离具有全局语义，等价于系统在故障下没有
良定义的行为；若缺少恢复规则，则系统在首次故障后不再回到规范状态。

综上，任一职能缺失都会使"良定义"这一前提被破坏。\qed
\end{proof}

\begin{remark}[定理的性质]\label{rmk:six-nature}
定理 \ref{th:six} 是一个\emph{条件必然性}结果，而非无条件定理。它的条件是
$|E| > 1$、$\mathcal{R}$ 有限且受争用、以及故障可达。它的证明是结构性的枚举论证，
而非某个形式系统中的推导。它可被证伪的方式是：给出一个满足前提却缺少六项之一的
良定义系统。该判断本身可被挑战。
\end{remark}

\begin{remark}[逃逸通道，以及如何穿过它做检验]\label{rmk:six-escape}
"或者在该系统自身的代码中，或者在其宿主中"这一措辞，既使该定理可用，也使它容易被
不诚实地辩护。粗心读来它会变得不可证伪：任何缺少六项之一的系统都被说成是从其宿主
那里获得了该职能，而宿主总在某处，于是该主张不可能失败。这一措辞不是一个漏洞
——定理的内容是该职能\emph{在语义中的某处被定义}，这是关于语义的主张，而不是关于
文件所有权的主张——但它确实转移了反驳的负担，因此值得说明这一负担如今落在何处。

一次反驳必须给出一个系统与一个职能，使得 (i) 操作语义在定理 \ref{th:six} 的意义下
良定义，(ii) 前提成立，且 (iii) 该职能\emph{既从该系统、也从它运行的每一个宿主中}
缺席。第 (iii) 条是难以满足的那一条，而它之所以难，原因是有信息量的而非回避性的：
有名字的宿主在相关职能上是稠密的。操作系统定义了调度、分配与隔离；语言运行时定义了
身份与生命周期，且往往还有通信；浏览器则定义了全部六项。因此一次反驳必须指出一个其
\emph{整个}执行环境都缺少该职能的系统，这在实践上意味着：要么是一个完全没有宿主的
系统——裸机加静态分配，而这已被定理 \ref{th:six} 的前提排除——要么是一个其宿主已被
证明缺少该职能的系统，而后者是关于那个宿主的、可核查的主张。

由此得出两条关于该主张应当如何使用的后果。第一，该定理最好被读作一份\emph{清单式
义务}，而不是一项关于世界的发现：对六项中的每一项，说明它在何处被定义，若答案是
"宿主"，则说明是哪个宿主、依据什么证据。一份填不满某一行的层级清单并不因此被反驳；
它只是因此不完整，而这是一个可用的状态。第二，该定理\emph{并未}确立的三条性质应当
被分离出来、当作假设处理，而不是折进主张之中，因为它们是可独立检验、也可独立置疑的：
\begin{enumerate}[label=(\alph*)]
\item \term{充分性}：六项已经足够，即语义良定义不需要六项之外的任何东西。这正是
      猜想 \ref{conj:relevant} 从粗粒化一侧处理、并由假设 $H_{10}$ 检验的主张。
\item \term{独立性}：六项之中没有哪两项可以用其余各项定义。该清单的动机来自彼此不同
      的语义义务——可区分性、有限性与争用、就绪集合的基数、交互的宾语、复合的方式，
      以及故障的传播范围——但义务的彼此不同并不等于实现的彼此独立，系统经常同时实现
      其中若干项（一张能力表同时是命名、分配与隔离）。
\item \term{最小性}：六项无法再被约简。论证表明去掉任意一项都会在前提之下破坏良定义；
      它并未表明六项在所有分解中最小，只表明这一分解相对于其自身前提是冗余自由的。
\end{enumerate}
消融实验是检验这三条的自然工具：对每个职能，把它从一个满足前提的系统中移除，观察
语义是否真的变得不良定义，以及余下的职能能否吸收它的义务。该实验尚未被做过。
\end{remark}

定理 \ref{th:six} 的措辞"或者在该系统自身的代码中，或者在其宿主中"，是关键性的。它
把"职能是否必然出现"与"职能是否必须自建"两个问题分离开来。

\section{自建与外借}\label{sec:buildadopt}

\begin{corollary}[自建/外借二分]\label{cor:buildadopt}
由定理 \ref{th:six}，六职能必被定义，但定义的位置是一个设计选择。该选择由不变量
的\emph{领域关键性}决定：
\begin{enumerate}[label=(\alph*)]
\item 若某职能所承载的不变量是\emph{领域关键}的，即该不变量必须由本系统拥有的保证
      才能成立，则该职能必须自建；
\item 若某职能所承载的不变量是\emph{通用}的，即通用宿主能够提供强度足够的保证，
      则应当外借；
\item 对通用职能自建实现，将产生本文称为\emph{通用件重造}的失败模式：以更高的长期
      理解与维护成本，重新实现一个更差的通用机制。
\end{enumerate}
\end{corollary}

这一判据解释了"应用"与"平台"的边界为什么正好落在它现在的位置：
\emph{你被迫自己拥有的不变量集合，就是你内核的边界}。

\begin{example}\label{ex:buildadopt}
一个游戏引擎必须自建实体-组件存储与帧调度，因为通用宿主无法提供"亚毫秒确定性帧
预算"这一保证。一个数据库必须自建存储引擎与预写日志，因为通用文件系统无法提供
"崩溃后恢复到确定状态"这一保证。反之，一个 Web 应用应当外借垃圾回收、线程调度与
内存分配，因为这些职能的不变量是通用的，自建只会得到更差的版本。
\end{example}

\begin{remark}[与企业理论的平行，以及它在哪里停止]\label{rmk:coase}
推论 \ref{cor:buildadopt} 回答的问题，在软件之外有一个发展成熟的对应物：一个组织应当
\emph{自建}什么，又应当\emph{外购}什么？科斯的回答是：企业的边界划在"在市场中交易的
成本超过在内部组织同一活动的成本"之处\cite{coase1937}，而后续关于自建/外购决策的
文献把这一点转化为治理成本的可操作比较\cite{williamson1985}。这一结构上的平行足够
接近，值得明确说出：在两种情形中，边界都不是关于工作的技术事实，而是哪些成本落在
边界的哪一侧的后果。

这里的判据在一个方面更锐利、在另一个方面更狭窄，而这一差别正是引用该平行而非依赖
它的理由。更锐利：科斯的比较是在两种成本之间，而这两种成本出了名地难以测量；而
推论 \ref{cor:buildadopt} 先问一个答案确定的定性问题——该不变量是否是只有本系统才能
拥有的那一个？——然后才考虑成本。一个拥有领域关键不变量的企业没有自建/外购决策可做，
无论市场替代品多么便宜。更狭窄：该推论没有涉及若干系统各自拥有某不变量的一部分、
或所有权可以按契约分割的情形，而那正是经济学文献大部分内容所在之处。

这一平行带来的不是结论，而是一个警告。如果软件边界像企业边界那样由治理成本划定，
那么它会表现出同样的病理：在理由消失之后仍然存续的边界；出于控制而非成本的理由被
选中的集成；以及地位受边界自身的转换成本、而非所供之物质量保护的供应商。这些在软件
中是可以观察到的，而一个声称解释该边界的框架应当能预测它们。本文没有做到这一点，这
是一个缺口，而不是一项主张。
\end{remark}

\begin{remark}[一个更锐利的判据：单调性]\label{rmk:calm}
推论 \ref{cor:buildadopt} 用不变量是否领域关键来判定自建还是外借。对"协调"这一类
职能，存在一个更锐利且可静态检查的判据：\emph{计算是否单调}。在单调性成立时\cite{calm2011}
（输出随输入增长而增长，且不依赖否定或聚合的全局知识），协调不是必需的，可以下推
到各参与者自行完成；非单调性（唯一性约束、否定、阈值聚合、非交换更新）则强制某个
位置承担协调。

这把第 \ref{sec:six} 节的职能 (F5)（通信与同步）与 (F3)（调度）的一部分从"规模到了
就必须集中"精确化为"非单调就必须集中"。对单调子系统，集成成本定理
（定理 \ref{th:integration}）中的 $k$ 可以被降到接近常数，因为参与者之间不需要就
共享状态达成一致。
\end{remark}

\section{跨领域对照}\label{sec:instances}

定理 \ref{th:six} 预测：横看每一列是一家不同公司、不同语言、不同时代、不同领域的
系统；纵看每一行是同一个器官。表 \ref{tab:organs} 给出对照。该表不主张这些机制可以
一一等同，而是把它们放在相同的\emph{管理职责}维度上比较。

\begin{table}[htbp]
\centering
\scriptsize
\setlength{\tabcolsep}{3pt}
\caption{同一组管理职责在不同领域的独立实现}\label{tab:organs}
\begin{tabular}{L{1.4cm}L{2.1cm}L{2.15cm}L{1.95cm}L{1.9cm}L{2.15cm}L{2.15cm}}
\toprule
\textbf{职能} & \textbf{Linux} & \textbf{Chromium} & \textbf{Tokio} & \textbf{Bevy}
& \textbf{Unreal} & \textbf{PostgreSQL} \\
\midrule
身份/Lifecycle & task\_struct, inode, file & Frame, Document, Node &
Task, JoinHandle & Entity, Component & UObject, Actor & backend, xid, tuple \\
分配/回收 & buddy, slab, page cache & PartitionAlloc, Oilpan &
slab 式任务分配 & archetype 列存 & FMemory, GC & MemoryContext, shared\_buffers \\
调度 & EEVDF, workqueue & task queue, RendererScheduler &
work-stealing, waker & Schedule/System 并行执行器 & TaskGraph, GameThread &
postmaster, autovacuum \\
命名/寻址 & kobject, sysfs, fd & DOM, Web IDL & Handle, waker &
Entity ID & FName, 反射元数据 & OID, 系统目录 \\
通信/同步 & syscall, netlink, IPC & Mojo IPC & mpsc, oneshot &
Event, Resource & 委托, 复制 & libpq, LWLock, 扩展钩子 \\
隔离/恢复 & namespace, cgroup, OOM & site isolation, 渲染进程沙箱 &
panic 传播, 取消 & 状态机隔离 & crash reporter & MVCC 快照, WAL 重放 \\
\bottomrule
\end{tabular}
\end{table}

反过来说，这六个器官在不同宿主中呈现的\emph{形态}差异极大，而这正是第
\ref{ch:convergence} 章"主干趋同、枝叶各异"的具体内容。形态差异的根源可以被系统
地列出：

\begin{enumerate}[label=(\alph*)]
\item \term{宿主内存模型}。手工内存管理的系统（C/C++/Rust）必须自建分配器甚至自建
      回收器（引擎对象回收、浏览器 Blink 的垃圾回收）；具备托管内存的系统把同样的
      精力花在停顿行为上，因为分配是现成的，痛点变成了延迟尾部分布。
\item \term{宿主并发模型}。运行时模型决定调度器官的形状：游戏引擎是帧驱动的协作式
      调度（可预测延迟优先于吞吐），服务器是事件驱动的抢占式调度（吞吐优先）。因此
      帧调度器与工作窃取调度器是同一个器官的两个亚种。
\item \term{领域原语粒度}。通用内核给出的是页、线程、文件描述符、事件轮询；但游戏
      引擎需要帧、组件与资产，数据库需要缓冲池、事务与版本链，浏览器需要文档节点、
      安全上下文与跨进程对象。原语粒度错配逼着每个领域自建基础设施。
\item \term{演化阶段}。这些内部操作系统本身也在持续被重构，且每次重构都是朝着
      "更完整的六职能"收敛而非发散。
\end{enumerate}

连操作系统自己也会被其用户绕过：高性能网络与存储栈在用户态轮询重做
I/O，异步 I/O 接口被重新设计，连操作系统都逃不过被自己的用户态"重新发明"的命运。
这不是反常，而是推论 \ref{cor:threshold} 的直接结果：当用户态的 $n \times m$ 越过
阈值，用户态就会长出新的内核。

\section{传输与重新导出}\label{sec:transport}

表 \ref{tab:organs} 显示同一组职能在不同领域中反复出现。但"反复出现"有两种截然不同的
解释，它们的预测完全不同：

\begin{enumerate}[label=(\alph*)]
\item \term{同构}：六职能是同一个结构的六次实现，可以彼此翻译；
\item \term{沿基底的传输}：六职能是某个更一般结构在不同\emph{问题形状}上的纤维，
      形状改变时结构必须重新导出，不能搬运。
\end{enumerate}

原文只讲了 (a)。区分二者需要显式引入"问题形状"这个基底。

\begin{definition}[问题形状与层族]\label{def:fibration}
称一个二元组 $\pi : \mathcal{L} \to \mathcal{P}$ 为一个\term{层族}，其中
\begin{itemize}
\item $\mathcal{P}$ 是\term{问题形状}的集合：由该领域的资源形态、并发结构、
      失效模式与时间尺度决定；
\item $\mathcal{L}$ 是层实现的集合；
\item 纤维 $\pi^{-1}(p)$ 是在问题形状 $p$ 下可行的层实现族。
\end{itemize}
$\pi$ 把"某个实现"映到"它解决的是什么形状的问题"。$\pi$ 的重索引给出
\term{沿基底的传输}：给定 $p \to p'$，把 $p$ 上的实现搬到 $p'$ 上。
\end{definition}

\begin{proposition}[传输的保持范围]\label{prop:transport}
设 $p \to p'$ 是问题形状之间的映射。若该映射是\emph{等价}（相同的资源形态、并发
结构、失效模式与时间尺度），则 $\mathcal{K}_6$ 的\emph{槽位}（每个职能所必须回答的
问题）沿传输保持，实现形态可以改变但不增加或减少槽位。若该映射不是等价，则至少
一个槽位的语义要求发生改变，实现必须重新导出。
\end{proposition}

\begin{proof}[论证]
$\mathcal{K}_6$ 的每一项由定理 \ref{th:six} 的语义义务唯一确定：身份由实体的可区分性
确定，分配由资源的有限性与争用确定，调度由就绪集合的基数确定，命名由交互的宾语
确定，通信由并发转移的复合方式确定，隔离由故障的传播范围确定。这六项义务只依赖
问题形状的四个分量。当映射保持这四个分量时，义务集合不变；当任一分量改变时，
对应义务的语义要求随之改变。
\end{proof}

\begin{corollary}[一个可观测的判别]\label{cor:transport-test}
由命题 \ref{prop:transport}，形态差异可以分成两类：\emph{同一槽位在不同基底上的
实现差异}（例如帧驱动与事件驱动的调度器，基底的时间尺度不同），与
\emph{不同槽位的差异}。前者不构成反例，后者构成。
\end{corollary}

推论 \ref{cor:transport-test} 使第 \ref{ch:convergence} 章的中命题（假设
\ref{hyp:medium}）可以被更严格地检验：比较两个系统时，必须先在基底上对齐，
再比较槽位。在基底上不对齐而直接比较模块，是既有架构比较研究最常犯的错误。

\begin{remark}[这一节的地位]\label{rmk:fibration-status}
定义 \ref{def:fibration} 与命题 \ref{prop:transport} 是一份\emph{结构提案}，不是
关于范畴论的既有定理的应用。它们的作用是把"同构还是传输"这个此前未被提出的问题
显式化，并给出一个可判别的推论（推论 \ref{cor:transport-test}）。把 $\pi$ 提升为
严格的纤维化并证明重索引函子的性质，是尚未完成的工作（第 \ref{ch:objections} 章
O 项）。
\end{remark}

\section{粗粒化判据}\label{sec:coarsegrain}

定理 \ref{th:six} 论证了六职能必然\emph{存在}，但没有说明它们是否是最\emph{显著}的
结构。一个更强的命题是：六职能是系统结构在粗粒化下的\emph{存活集}。

\begin{definition}[粗粒化序列]\label{def:coarsegrain}
设系统的架构关系型为 $\mathcal{G}$（定义 \ref{def:relgraph}）。一个\term{粗粒化}
是一个商映射 $q : \mathcal{G} \to q(\mathcal{G})$，它合并构件并保留关系类型。
称一列商映射 $q_1, q_2, \dots$（每一步的粒度严格变粗）为一个\term{粗粒化序列}。
结构 $\mathcal{S}$ 在该序列下\term{存活}，若存在 $N$ 使得对所有 $n \ge N$，
$\mathcal{S}$ 在 $q_n(\mathcal{G})$ 中仍可识别。
\end{definition}

\begin{conjecture}[六职能是相关结构]\label{conj:relevant}
在一个正常偏置的大型系统的粗粒化序列下，存活的结构恰为 $\mathcal{K}_6$ 的槽位及
其关系型；其余结构（具体模块、算法选择、数据布局细节）在足够粗的粒度上被冲掉。
\end{conjecture}

猜想 \ref{conj:relevant} 的价值在于它给出一个\emph{单项目内}的证伪程序：不需要跨项目
统计比较，只需对一个代码库反复粗粒化，观察存活集。若存活集不是六职能，猜想被否定，
且第 \ref{ch:convergence} 章的收敛论题需要一个不同的解释。

\begin{principle}[实施形式]\label{pr:coarsegrain-impl}
粗粒化可以完全机械地实施：以模块依赖图为基础，反复施加三类操作：(i) 合并扇入扇出
低于阈值的节点；(ii) 折叠可传递的依赖边；(iii) 抹去不携带不变量（定义
\ref{def:enforce} 中 $\mathrm{S}(\iota) = \varnothing$）的节点。每一步记录存活结构的
变化。整个实验只需一个仓库与一个依赖图，不需要标签、不需要人工分类，因此不受
第 \ref{sec:limits} 节指出的"职责分类系统选择自由度"影响。
\end{principle}

第 \ref{ch:program} 章把猜想 \ref{conj:relevant} 写成假设 $H_{10}$。

\begin{remark}[与重整化的关系]\label{rmk:rg}
原则 \ref{pr:coarsegrain-impl} 的操作与重整化群的粗粒化在形式上相同：消去短程自由度，
观察长程结构是否收敛到不动点。这里的不动点就是六职能槽位。必须明确的是，这
\emph{不}意味着软件具有重整化群流，因为软件没有尺度分离（原则 \ref{pr:vm} 之前的讨论），
因而不存在真正意义上无关算符的指数衰减。相同的是\emph{操作}，不是\emph{机制}。
\end{remark}

\section{虚拟机化与深模块}\label{sec:virtualisation}

"压缩复杂度、为上层提供高级语言与符号"这一表述可以精确化。每个基础设施层的核心
工作都是同一种动作：

\begin{principle}[虚拟机化]\label{pr:vm}
基础设施层把\emph{物理的、连续的、易变的}东西，翻译成\emph{离散的、稳定的、
可组合的}领域级符号。
\end{principle}

\begin{example}\label{ex:vm}
操作系统：物理内存 $\to$ 虚拟地址空间；块设备 $\to$ 文件；硬件中断 $\to$ 信号。
数据库：物理页 $\to$ 多版本元组 $\to$ SQL 中的"行"。
浏览器：字节流 $\to$ 文档树；套接字事件 $\to$ 事件循环回调。
游戏引擎：内存中的裸字节 $\to$ 组件与资源；一帧的时间 $\to$ 调度表。
异步运行时：事件轮询位图 $\to$ \code{async}/\code{await} 中的 Future。
\end{example}

用模块设计的术语\cite{parnas1972,ousterhout2018}，这些全是\emph{深模块}：
接口极小，实现极深。一行分配调用背后是
整个分配器体系；一行任务创建背后是整个调度器。基础设施层的职责就是持续生产深模块，
把复杂度封在下面，把干净的符号暴露在上面。

这一点把它与第 \ref{ch:comprehension} 章的理解理论直接连接起来：深模块正是压缩率
（定义 \ref{def:compression}）极高的层级，它商掉了大量底层差异，同时保持了上层
所需的预测能力。这也解释了为什么好的基础设施同时是"高性能的"与"易理解的"：两者都
是同一个商映射的后果。

\section{智能体系统的职能分化}\label{sec:agent}

如果定理 \ref{th:six} 成立，那么最新出现的复杂计算域应当正在重演这一演化。表
\ref{tab:agent} 给出观察。

\begin{table}[htbp]
\centering
\small
\caption{智能体系统对六职能的重新发现}\label{tab:agent}
\begin{tabular}{L{3.4cm}L{11.0cm}}
\toprule
\textbf{职能} & \textbf{在智能体域中的对应物} \\
\midrule
分配与回收 & 上下文窗口作为快速存储；检索增强作为缺页处理；外置记忆作为交换区；
上下文压缩作为内存紧缩 \\
通信与同步 & 工具调用或函数调用作为系统调用原语 \\
调度 & 智能体主循环、多智能体编排 \\
身份与生命周期 & 会话、任务与子目标的生命周期管理 \\
命名与寻址 & 工具注册与发现协议 \\
隔离与恢复 & 沙箱、重试、检查点与回滚 \\
\bottomrule
\end{tabular}
\end{table}

这一观察的价值在于它是一个\emph{自然实验}：没有人在刻意模仿内核，是复杂度本身把
同样的解剖结构压了出来。它同时构成假设 $H_3$ 的经验基础（第 \ref{ch:program} 章）。

定理 \ref{th:integration} 与 \ref{th:six} 在给定前提下成立，其可挑战之处在于前提与
结构论证的严密性。经验上可证伪的是：六职能的出现是否由 $n \times m$ 而非代码行数
预测（假设 $H_3$ 与 $H_{10}$）；以及是否存在满足前提却长期缺少某一职能的良定义系统。
